\chapter[2004 --- Axel et al.]{2004: Richard Axel, Linda B. Buck}
\label{ch:2004_Axel}

\section*{Main Contribution}

\textit{Discovered odorant receptors and the organization of the olfactory system, explaining how the sense of smell detects and discriminates thousands of different chemicals.}\cite{nobel-2004,lecture-2004}

\section{Official Citation}

for their discoveries of odorant receptors and the organization of the olfactory system

\section{Background and Historical Context}

The sense of smell---olfaction---is among the most ancient and fundamental of the senses. It enables animals to find food, avoid predators, locate mates, and navigate their environment. Humans can detect and discriminate thousands of different odorants, ranging from the floral scent of rose to the acrid smell of smoke. Yet for most of the twentieth century, the molecular mechanisms underlying olfaction remained among the least understood of all sensory systems. The discovery of odorant receptors and the organization of the olfactory system by Richard Axel and Linda B. Buck, recognized with the 2004 Nobel Prize in Physiology or Medicine, provided the molecular and neural framework for understanding how the brain encodes and processes olfactory information.

Richard Axel was born on July 2, 1946, in Brooklyn, New York. He studied medicine at Johns Hopkins University, earning his M.D.\ in 1970, and then joined the faculty at Columbia University, where he became University Professor and investigator at the Howard Hughes Medical Institute. Axel's early research concerned the molecular biology of gene expression, but he became increasingly interested in the problem of olfaction in the 1980s.

Linda Buck was born on January 29, 1947, in Seattle, Washington. She studied psychology at the University of Washington and then earned her Ph.D.\ in immunology at the University of Texas Southwestern Medical Center in 1980 under the supervision of Howard Goodman. After postdoctoral work with Axel at Columbia University, Buck joined the faculty at Harvard Medical School and later moved to the Fred Hutchinson Cancer Research Center in Seattle.

The intellectual context of their work was the recognition that the olfactory system faces a unique coding challenge. Unlike the visual system, which uses three types of photoreceptors to encode the entire spectrum of visible light, or the auditory system, which encodes sound as a spatial map of frequencies on the basilar membrane, the olfactory system must distinguish among thousands of different chemical structures---and must do so using a combinatorial coding strategy in which each odorant is represented by a unique pattern of neural activity across many different receptor types. The molecular basis of this combinatorial code was unknown until Axel and Buck made their discoveries.

\section{The Discovery/Contribution}

\subsection{The Discovery of Odorant Receptor Genes}

In 1991, Axel and Buck published a landmark paper in \textit{Cell} in which they reported the identification of a large family of genes encoding odorant receptors---G protein--coupled receptors (GPCRs) expressed in the olfactory epithelium of the nose. Their approach was based on a simple but powerful strategy: they reasoned that olfactory receptor neurons must express receptor proteins that bind odorant molecules, and that these receptors should be detectable as a family of related genes expressed specifically in the olfactory epithelium.

Using a combination of molecular cloning and hybridization techniques, Axel and Buck identified 13 related genes in the rat olfactory epithelium that encoded putative seven-transmembrane-domain proteins---the characteristic structure of GPCRs. These proteins, which they named odorant receptors (ORs), were expressed exclusively in the olfactory epithelium and were not found in any other tissue. The discovery was remarkable because the 13 genes they identified represented only a fraction of the olfactory receptor gene family---later work revealed that mice have approximately 1,100 odorant receptor genes, and humans have approximately 400.

\subsection{The Combinatorial Code}

The identification of odorant receptors immediately raised the question: how do a few hundred receptor types encode the discrimination of thousands of different odorants? Axel and Buck proposed the combinatorial coding hypothesis, which states that each odorant activates a unique combination of receptor types, and each receptor type can be activated by multiple odorants. The identity of an odorant is encoded not by the activity of a single receptor but by the pattern of activity across many receptors---a ``combinatorial code'' analogous to the way letters are combined to form words.

This hypothesis was supported by subsequent experimental studies. Using calcium imaging and electrophysiological techniques, Axel and Buck and their colleagues showed that individual odorant receptor neurons express only one receptor type (the ``one neuron--one receptor'' rule) and that individual odorants activate characteristic combinations of receptor neurons. The pattern of activated receptors is then transmitted to the olfactory bulb, where it is processed by glomeruli---structures that receive input from receptor neurons expressing the same receptor type.

\subsection{The Olfactory Map and Topographic Organization}

Buck's subsequent work revealed the neural circuitry by which olfactory information is processed in the brain. She showed that axons from olfactory receptor neurons expressing the same receptor converge on one or two specific glomeruli in the olfactory bulb, creating a topographic ``odor map'' in which the spatial pattern of glomerular activation reflects the pattern of receptor activation. This map is then transmitted to higher brain regions---the piriform cortex, the amygdala, and the entorhinal cortex---where it is decoded and associated with memories, emotions, and behaviors.

Buck also discovered that the olfactory system exhibits remarkable plasticity: when an odorant receptor is experimentally deleted, the affected receptor neurons lose their identity and can adopt the identity of other receptor types, suggesting that the maintenance of the olfactory map requires ongoing receptor expression.

\subsection{Vomeronasal and Pheromone Receptors}

Buck and her colleagues also identified receptors in the vomeronasal organ---a specialized olfactory structure in many animals that detects pheromones, the chemical signals used for intraspecies communication. These receptors, which belong to a distinct gene family from the odorant receptors, mediate innate behavioral responses to pheromones, including mating behavior and aggression. The identification of vomeronasal receptors extended the molecular framework of olfaction beyond odor detection to the broader realm of chemical communication.

\section{Impact on Modern Medicine}

The discoveries of Axel and Buck have had a significant impact on neuroscience, medicine, and our understanding of sensory perception.

\subsection{Olfactory Disorders}

The identification of odorant receptors has provided the molecular basis for understanding olfactory disorders, including anosmia (loss of smell) and hyposmia (reduced sense of smell). These conditions affect millions of people worldwide and can result from viral infections (including COVID-19), head trauma, neurodegenerative diseases, and aging. Understanding the molecular targets of olfactory loss has guided the development of diagnostic tests and potential therapies for olfactory dysfunction.

\subsection{Neurodegenerative Disease}

Olfactory dysfunction is one of the earliest symptoms of several neurodegenerative diseases, including Parkinson's disease and Alzheimer's disease. In Parkinson's disease, the loss of the sense of smell often precedes motor symptoms by years or decades. Understanding the mechanisms of olfactory dysfunction in these diseases may provide the basis for early diagnostic biomarkers and preventive interventions.

\subsection{Chemical Ecology and Pest Control}

The identification of insect odorant receptors has opened new avenues for pest control. Insects rely heavily on olfaction for finding hosts, mates, and oviposition sites. The development of odorant receptor--based traps and repellents---targeting mosquito odorant receptors, for example---represents a new approach to combating vector-borne diseases such as malaria, dengue, and Zika.

\subsection{Olfactory-Based Diagnostics}

The olfactory system is exquisitely sensitive to certain volatile organic compounds that are biomarkers of disease. The detection of specific odorants in breath or urine has been explored as a non-invasive diagnostic tool for lung cancer, diabetes, and other conditions. Understanding the molecular specificity of odorant receptors has informed the development of electronic noses---sensor arrays that mimic the olfactory system's ability to detect and discriminate volatile chemicals.

\section{Legacy}

Richard Axel continued his work at Columbia University, where he is University Professor and investigator at the Howard Hughes Medical Institute. He has made additional contributions to the understanding of olfactory coding and the neural circuits underlying olfactory behavior. He has received numerous honors, including the Gairdner Foundation International Award.

Linda Buck continued her work at the Fred Hutchinson Cancer Research Center, where she became a member of the Basic Sciences Division and an affiliate professor at the University of Washington. She received the Richard Lounsbery Award and numerous other honors. Buck continues to study the neural mechanisms of olfactory perception and the genetic basis of individual differences in smell.

Together, Axel and Buck solved one of the great mysteries of sensory biology: how the nose detects and discriminates the vast universe of odorant molecules. Their discovery of the odorant receptor gene family and the combinatorial code for olfactory perception has provided the framework for understanding one of the most complex and evolutionarily ancient sensory systems, with implications that extend from basic neuroscience to clinical medicine and biotechnology.


\section{Modern Developments}

Axel and Buck's olfactory work has led to modern sensory neuroscience. Understanding of odorant receptors has informed development of artificial noses for disease diagnosis. Understanding of olfactory circuitry has informed development of treatments for anosmia (loss of smell), a common symptom of COVID-19. Understanding of olfactory coding has inspired AI-based pattern recognition systems. Olfactory dysfunction is now recognized as an early biomarker for Alzheimer's and Parkinson's diseases.


\section{Bibliography}

\begin{thebibliography}{9}
\bibitem{nobel-2004}
Nobel Prize Outreach, ``The Nobel Prize in Physiology or Medicine 2004,''\emph{NobelPrize.org}.\\
\url{https://www.nobelprize.org/prizes/medicine/2004/summary/}

\bibitem{lecture-2004}
Richard Axel, ``Nobel Lecture,''\emph{NobelPrize.org}. Winner-authored primary source; downloadable from the official Nobel Prize archive.\\
\url{https://www.nobelprize.org/prizes/medicine/2004/axel/lecture/}
\end{thebibliography}
